\section{Jacobian 与多元链式法则}
\label{sec:02-Calculus/06-Multivariable-Calculus/04}

一个前馈网络就是一串向量到向量的映射接在一起：输入 \(\mathbf x\in\R^{d_0}\) 经过第一层变成 \(\R^{d_1}\) 里的向量，再变成 \(\R^{d_2}\) 里的，一路走到最后输出一个标量损失。训练要的是损失对第一层参数的导数。中间隔着几十层，每层的输出都是几百上千维，这个导数凭什么算得出来？

答案的全部数学内容就在本节。它不需要新概念，只需要把上一节那句话认真对待：\emph{导数是线性映射}。线性映射天然可以复合，复合的结果还是线性映射——链式法则说的就是这件事，剩下的都是记账。

\subsection{向量值函数的导数是一个矩阵}

设 \(F:\R^n\to\R^m\)，\(F=(F_1,\dots,F_m)\)，在 \(\mathbf a\) 可微。全导数 \(DF(\mathbf a)\) 是从 \(\R^n\) 到 \(\R^m\) 的线性映射，在标准基下写成 \(m\times n\) 矩阵，称为 Jacobian：
\[
J_F=
\begin{bmatrix}
\partial F_1/\partial x_1&\cdots&\partial F_1/\partial x_n\\
\vdots&&\vdots\\
\partial F_m/\partial x_1&\cdots&\partial F_m/\partial x_n
\end{bmatrix},
\qquad
F(\mathbf a+\mathbf h)\approx F(\mathbf a)+J_F(\mathbf a)\mathbf h .
\]
第 \(j\) 列回答``只动第 \(j\) 个输入，全部输出各变多少''，第 \(i\) 行是第 \(i\) 个输出分量的梯度转置。标量函数是 \(m=1\) 的特例，此时 \(J_f=\nabla f^{\trans}\)——梯度是列向量，Jacobian 是行向量，两者差一个转置，这个差别在拼接公式时必须盯住。

极坐标映射 \(F(r,\theta)=(r\cos\theta,\,r\sin\theta)\) 的 Jacobian
\[
J_F=
\begin{bmatrix}
\cos\theta&-r\sin\theta\\
\sin\theta&r\cos\theta
\end{bmatrix}
\]
按列读很有意思：第一列是径向的单位位移，第二列是角度变化引起的切向位移，长度带着因子 \(r\)——同样转过一个小角度，离原点越远走得越长。这个矩阵在\hyperref[sec:02-Calculus/06-Multivariable-Calculus/07]{``多重积分与变量替换''}一节还会出现，那里取的是它的行列式。

\subsection{复合的导数是映射的复合}

\begin{theorem}
若 \(F:\R^n\to\R^m\) 在 \(\mathbf x\) 可微，\(G:\R^m\to\R^p\) 在 \(F(\mathbf x)\) 可微，则 \(G\circ F\) 在 \(\mathbf x\) 可微，且
\[
D(G\circ F)(\mathbf x)=DG(F(\mathbf x))\circ DF(\mathbf x),
\qquad\text{即}\qquad
J_{G\circ F}(\mathbf x)=J_G(F(\mathbf x))\,J_F(\mathbf x).
\]
\end{theorem}

证明的骨架只有两步。把 \(F\) 在 \(\mathbf x\) 处的增量写成 \(J_F\mathbf h+\mathbf r_F\)，代进 \(G\) 的线性近似，得到
\[
G(F(\mathbf x+\mathbf h))-G(F(\mathbf x))
=J_G\bigl(J_F\mathbf h+\mathbf r_F\bigr)+\mathbf r_G .
\]
线性映射把高阶小量仍送成高阶小量（\(\norm{J_G\mathbf r_F}\le\norm{J_G}\norm{\mathbf r_F}\)），而 \(\mathbf r_G\) 相对于它的自变量 \(J_F\mathbf h+\mathbf r_F\) 高阶，后者又与 \(\norm{\mathbf h}\) 同阶或更小。两项合并仍是 \(o(\norm{\mathbf h})\)，剩下的就是 \(J_GJ_F\mathbf h\)。真正做功的是可微性的``统一''那一层：误差用同一个 \(\norm{\mathbf h}\) 度量，才允许这样一层层往外传。

值得强调的是等式左边那句话——复合映射的导数就是导数的复合。矩阵乘法在这里没有任何独立地位，它只是线性映射复合在选定基下的记账方式。写成分量，
\[
\frac{\partial z}{\partial x_j}=\sum_{i}\frac{\partial z}{\partial u_i}\frac{\partial u_i}{\partial x_j},
\]
读作``所有中间路径的贡献相加''；写成矩阵，读作``两个线性映射串起来''。后者不但更短，还保留了尺寸信息：\(p\times m\) 乘 \(m\times n\) 得 \(p\times n\)，顺序颠倒立刻尺寸不合，这是一道免费的校验。其实链式法则在日常生活中每天都在发生。想知道车每小时烧多少油，你会把``每小时跑 \(80\) 公里''和``每公里耗油 \(0.08\) 升''乘起来，得到每小时 \(6.4\) 升。这里的``公里''是中间量，算完就约掉了，正如矩阵乘法里中间那个维数 \(m\)；要是乘完单位对不上，多半是哪一环接错了。

\subsection{同一个乘积，两种算法}

矩阵乘法满足结合律，于是链条
\[
J=J_L\,J_{L-1}\cdots J_2\,J_1
\]
可以从任一端开始算。这不是排版偏好，它是自动微分两种模式的全部区别。

\begin{center}
\begin{tikzpicture}[x=1cm,y=1cm,font=\small,>=stealth]
  \draw[black!55,fill=blue!6] (0,0) ellipse (0.85 and 1.15);
  \draw[black!55,fill=blue!6] (3.4,0) ellipse (0.85 and 1.15);
  \draw[black!55,fill=blue!6] (6.8,0) ellipse (0.85 and 1.15);
  \node at (0,0) {\(\R^{n}\)};
  \node at (3.4,0) {\(\R^{m}\)};
  \node at (6.8,0) {\(\R^{p}\)};
  \draw[->,thick] (0.95,0.25) -- (2.45,0.25);
  \draw[->,thick] (4.35,0.25) -- (5.85,0.25);
  \node[above] at (1.7,0.3) {\(J_F\)  \((m\times n)\)};
  \node[above] at (5.1,0.3) {\(J_G\)  \((p\times m)\)};
  \draw[->,black!60,thick] (0.3,-1.25) .. controls (2.5,-2.3) and (4.7,-2.3) .. (6.5,-1.25);
  \node[black!70] at (3.4,-2.48) {\(J_{G\circ F}=J_G\,J_F\)};
  \node[right,align=left] at (-0.5,-3.7)
    {前向模式：\(J_G(J_F\mathbf v)\)，从输入侧推一个方向向前走；\\
     反向模式：\((\mathbf u^{\trans}J_G)J_F\)，从输出侧拉一个权重向后走。};
\end{tikzpicture}

\emph{括号加在哪一边，决定了每一步是矩阵乘向量还是矩阵乘矩阵——两者的代价可以差好几个数量级。}
\end{center}

假设各层宽度都是 \(d\)，共 \(L\) 层，最终输出是一个标量损失。若老老实实先把 \(J\) 这个矩阵乘出来，每次矩阵乘矩阵要 \(O(d^3)\)，全程 \(O(Ld^3)\)，而且中间要存下 \(d\times d\) 的稠密矩阵。反向模式从末端出发：先拿标量输出对应的 \(1\times d\) 行向量，右乘一个 Jacobian 得到新的 \(1\times d\) 行向量，如此一路倒推回输入端。每步是向量乘矩阵，\(O(d^2)\)，全程 \(O(Ld^2)\)。省下的那个 \(d\) 倍，正是训练大模型在计算上可行的原因。

代价不对称的根子在维度：反向模式的优势来自输出只有一维。这也解释了一条工程惯例——损失函数必须是标量。若要对一个 \(k\) 维输出求完整 Jacobian，反向模式得跑 \(k\) 遍，此时若输入维度更小，前向模式反而划算。框架里真正被反复调用的原语从来不是``那个矩阵''，而是向量与 Jacobian 的乘积；把它当成一个可以直接实现、不必显式构造矩阵的操作，是\hyperref[sec:10-Matrix-Calculus/04-Applications/03]{``计算图上的反向传播与自动微分''}一节的出发点。

\begin{example}
令 \(F(x,y)=(u,v)=(x+y,\,xy)\)，\(g(u,v)=u^2+e^{v}\)。则
\[
J_F=\begin{bmatrix}1&1\\ y&x\end{bmatrix},
\qquad
J_g=\begin{bmatrix}2u&e^{v}\end{bmatrix},
\]
于是
\[
J_{g\circ F}
=\begin{bmatrix}2(x+y)&e^{xy}\end{bmatrix}
\begin{bmatrix}1&1\\ y&x\end{bmatrix}
=\begin{bmatrix}2(x+y)+ye^{xy}&\;2(x+y)+xe^{xy}\end{bmatrix}.
\]
把 \(g\circ F\) 展开成 \((x+y)^2+e^{xy}\) 再逐个求偏导，结果当然一样，但那样做会把中间变量的结构碾平；层数一多就不再可行。
\end{example}

\subsection{记号约定要在动笔前说清}

有的资料把梯度写成行向量，有的写成列向量；有的按``分子布局''排导数，有的按``分母布局''。两套都自洽，混用则必错。本书统一用梯度列向量、Jacobian 的输出分量按行排，因此 \(J_f=\nabla f^{\trans}\)。读外部材料时先看一眼向量是行是列、矩阵是几乘几，再决定公式能不能直接接过来。\hyperref[sec:10-Matrix-Calculus/01-Derivative-Conventions/02]{``形状先于公式，布局必须声明''}一节把这件事做成了固定的检查动作。

\subsection{行列式非零意味着没有方向被压扁}

到目前为止，Jacobian 扮演的角色是复合时的记账工具。但既然它就是局部线性化的全部信息，一个自然的问题浮上来：这个线性模型什么时候没有把信息弄丢——什么时候映射在局部能被反过来读，输出能唯一地还原出输入？当 \(F:\R^n\to\R^n\) 时，Jacobian 是方阵，这个问题由它的行列式回答。

\begin{theorem}[逆函数定理]
若 \(F\) 在 \(\mathbf a\) 附近连续可微且 \(\det J_F(\mathbf a)\ne0\)，则 \(F\) 在 \(\mathbf a\) 的某个邻域上有可微的局部逆，且
\[
J_{F^{-1}}(F(\mathbf a))=J_F(\mathbf a)^{-1}.
\]
\end{theorem}

几何上，\(\det J_F\ne0\) 说的是局部线性化没有把任何方向压成零——小球被拉成一个不退化的椭球。注意这是纯粹的局部结论：极坐标映射处处（除原点外）行列式非零，却因为角度以 \(2\pi\) 为周期而远非全局单射。

行列式恰好为零很少发生，接近零才是常态，而后果是连续的。机械臂的关节角 \(\mathbf q\) 决定末端位置 \(F(\mathbf q)\)，速度关系是 \(\dot{\mathbf x}=J_F(\mathbf q)\dot{\mathbf q}\)。在奇异姿态附近，某个末端运动方向对应的关节速度需求趋于无穷——不是公式里出现了除零，而是 Jacobian 的某个奇异值趋于零。同样的量在数值线性代数里叫条件数，在优化里表现为病态的搜索方向，在深度网络里表现为梯度沿某些方向被逐层放大或压缩。真正该盯的从来不是行列式这一个数，而是奇异值的分布。

一阶信息到此收束：偏导数只看坐标方向，Jacobian 把所有方向组织成一个线性映射，链式法则处理复合，行列式与奇异值报告这个映射有没有退化。下一节转向二阶——当梯度归零、一阶模型什么也不说时，是哪个矩阵接手回答``这里是坑、是包，还是马鞍''。
